\documentclass[10pt,a4paper]{amsart}
\usepackage[margin=19mm]{geometry}
\usepackage{amsmath,amssymb,mathtools}
\usepackage[scr=rsfs,cal=euler]{mathalfa}
\usepackage{xfrac}
\usepackage[unicode,hidelinks,bookmarks=false]{hyperref}
\newcommand{\ie}{i.e.\ }
\newcommand{\cf}{cf.\ }
\newcommand{\resp}{respectively\ }
\newcommand{\Subsection}{\subsection}
\providecommand{\nobreakdash}{\nobreak\hbox{-}}
\setlength{\emergencystretch}{2em}
\multlinegap0pt
\usepackage[normalem]{ulem}
\usepackage{dsfont}
\usepackage{bm}
\usepackage{amssymb}
\usepackage{xcolor}
\usepackage{tcolorbox}
\usepackage{tikz}
\newtheorem{lemma}{Lemma}
\newcommand{\bR}{\mathbb{R}}
\newcommand{\cO}{\mathcal{O}}
\title{Convergence of Stochastic Gradient Langevin Dynamics in the Lazy Training Regime}
\author{Noah Oberweis, Semih Cayci}
\date{Source: arXiv:2510.21245v3 (CC BY 4.0)}
\begin{document}
\maketitle

\noindent\emph{Source and licence.} Convergence of Stochastic Gradient Langevin Dynamics in the Lazy Training Regime. Version arXiv:2510.21245v3 (CC BY 4.0), distributed by arXiv under the Creative Commons Attribution 4.0 International licence, http://creativecommons.org/licenses/by/4.0/, as recorded in the arXiv licence metadata for this exact version and shown on the abstract page https://arxiv.org/abs/2510.21245v3. Full licence notice: \texttt{licenses/arxiv-2510.21245-CC-BY-4.0.txt}.\par\smallskip

\noindent\textit{Editorial scope.} Counted unit: the article's lemma lemma:curgvature (source0.tex lines 1236--1238). The article's complete original proof of it is delivered with it (source0.tex lines 1239--1275, one \texttt{proof} environment of 37 lines). No mathematics is written, repaired or re-derived: the statement and the proof are the article's own lines, in the article's own notation, and no retained line is substituted or re-flowed.  No further source-local context is needed: every cross-reference of every retained line resolves inside the two delivered spans.  Nothing was omitted from any retained span, so the package carries no omission record.  No retained line contains a citation, so the wrapper carries no bibliography and no bibliography entry.  No retained line contains a figure environment, an image-inclusion command or drawing code, so no image is delivered and the package declares no asset dependency.  Editorial formatting only, in addition: an AMS \texttt{amsart} preamble that restates the article's own declarations and macros used by the retained lines, namely the environments \texttt{lemma} on separate counters, together with the standard packages those lines need.  The retained environments print in this document as Lemma 1 in order of appearance, whereas the article prints the same items with the numbering of its own full document.  The complete native source of the article as distributed in the arXiv e-print for this exact version is bundled unchanged as \texttt{sources/187707/source0.tex}.
\begin{lemma}\label{lemma:curgvature}[Curvature of parameter-loss-mapping]
    There exist parameters $\alpha,\eta_\alpha>0$ such that $\nabla^2_\omega(R(\alpha h(\omega)))\preceq \frac{\alpha^2}{\eta_\alpha} Id$.
\end{lemma}

\begin{proof}
    To prove this property, we need to calculate the Hessian of the mapping $\omega \mapsto R(\alpha h(\omega))$ and then determine the maximum eigenvalue of the Hessian. To do this, we will first decompose the Hessian into different parts, and then write the eigenvalues of the components using block-submatrices. 
    The gradient of the above function is given by
    $$\nabla_\omega R(\alpha \phi(x;\omega,c))=\frac{\alpha}{n}\sum_{i=1}^nr_i(\omega)J_i(\omega),$$
    where 
    \begin{align*}
        r_i(\omega):&=\alpha \phi(x_i;\omega,c)-y_i\\
        J_i(\omega)&=\nabla_\omega \phi(x_i;\omega,c).
    \end{align*}
    Using this notation, we can write the Hessian as 
    $$\nabla^2_\omega R(\alpha \phi(x;\omega,c))=\frac{\alpha^2}{n}\sum_{i=1}^nJ_i(\omega)J_i(\omega)^T+\frac{\alpha}{n}\sum_{i=1}^nr_i(\omega)\nabla^2_\omega\phi(x;\omega,c),$$
    which implies that 
    $$\lambda_{max}\left(\nabla^2_\omega R(\alpha \phi(x;\omega,c))\right)\leq\frac{\alpha^2}{n}\sum_{i=1}^n\|J_i(\omega)\|^2+\frac{\alpha}{n}\sum_{i=1}^n|r_i(\omega)|\|\nabla^2_\omega\phi(x;\omega,c)\|.$$
    In other words, we only need to bound $\|J_i(\omega)\|,|r_i(\omega)|$ and $\|\nabla^2_\omega \phi(x_i;\omega,c)\|$ in order to show the proposition. 
    \begin{itemize}
        \item We start off by bounding $\|J_i(\omega)\|^2$:
        \begin{align*}
            \|J_i(\omega)\|^2&= \sum_{j=1}^m\|\frac{c_j}{\sqrt{m}}\sigma'(\omega_j^T x_i)x_i\|^2\\
            &=\frac{\|x_i\|^2}{m}\sum_{j=1}^mc_j^2\sigma'(\omega_j^Tx_i)^2\\
            &\leq \frac{\|x_i\|^2}{m}\|c\|_2^2
        \end{align*}
        where we used $\forall x\in \bR:|\sigma'(x)|\leq 1$.
        \item The bound of $r_i$ follows immediately from  $$|r_i(\omega)|=|\alpha\phi(x_i;\omega,c)-y_i|\leq \alpha|\frac{1}{\sqrt{m}}\sum_{j=1}^m c_j\sigma(\omega_j^Tx)|+|y_i|\leq \frac{\alpha}{\sqrt{m}}\|c\|_1+|y_i|.$$
        \item Lastly, we still need to bound $\|\nabla^2_\omega \phi(x_i;\omega,c)\|$. We will do this by considering the blocks of the Hessian $\nabla^2_{\omega_j,\omega_k}\phi(x_i;\omega,c)$. First, we observe for $j\neq k$ 
        $$\nabla^2_{\omega_j,\omega_k}\phi(x_i;\omega,c)=D_{\omega_k}\left(\frac{1}{\sqrt{m}}c_j\sigma(\omega_j^Tx)x\right)\overset{j\neq k}{=}0,$$
        which means that this Hessian is a block-diagonal matrix. In addition, it holds for $j=k$
        $$\nabla^2_{\omega_j,\omega_j}(\phi(x_i;\omega,c))=\frac{1}{\sqrt{m}}c_j\sigma''(\omega_j^Tx_i)x_ix^T_i.$$
        Taking the norm of this equation, we receive
        $$\lambda_{max}\left(\nabla^2_{\omega_j,\omega_j}(\phi(x_i;\omega,c))\right)\leq \frac{4}{3\sqrt{3m}}c_j\|x_i\|,$$
        where we used $\forall x\in \bR:\|\sigma''(x)\|\leq \frac{4}{3\sqrt{3}}$.
    \end{itemize}
    Putting all of the above together, we get the upper bound
    \begin{align*}
        \lambda_{max}(\nabla_\omega^2R(\alpha\phi(x;\omega,c)))\leq \frac{\alpha^2\|c\|_2^2\|x\|_2^2}{mn}+\frac{4\alpha\|c\|_\infty}{3n\sqrt{3m}}\left(\frac{\alpha\|c\|_1}{\sqrt{m}}+\|y\|_\infty\right)\|x\|^2_2,
    \end{align*}
    which is in $\cO(\alpha^2)$ for a given dataset and predictor class. This implies that for $\eta_\alpha$ sufficiently small, the above proposition holds.
    \end{proof}

\end{document}
